\documentclass[11pt,a4paper]{article}
\usepackage[top=1in,bottom=1in,left=1.1in,right=1.1in]{geometry}
\usepackage{mathptmx}
\usepackage{hyperref}
\pagestyle{empty}
\begin{document}

\begin{flushright}
Hussain Touhid Siddiquee\\
Department of EEE, Leading University, Sylhet-3112, Bangladesh\\
E-mail: \url{touhidsiddiqueeraj@gmail.com}\\
7 October 2026
\end{flushright}

\vspace{1em}
\noindent Dear Editor,

\vspace{1em}
\noindent\textbf{Submission of Regular Paper: ``Modeling electrical losses in CsPbI3 solar cells: stacked idealizations and a routine-quality envelope at 18.0\%''}\\
by Hussain Touhid Siddiquee and Md.\ Abdul Malek Fahim

\vspace{1em}
Please consider our manuscript for publication as a Regular Paper in \textit{Solar Energy}.

\textbf{What the paper does.} On one fully worked ITO/TiO2/CsPbI3/CBTS/Ni stack in SCAPS-1D, we (i) locate a permissive-box ceiling of 24.95\% by joint optimisation of thickness, doping, defects, gap, and \textit{both} transport-layer affinities; (ii) sample a frozen routine-quality prior volume (96-point Latin hypercube with published, measured-anchored ranges) giving a prior-predictive envelope of median 18.0\% with 5--95\% range 15.6--20.0\%, above which the 2024--2026 experimental records fall; and (iii) price every load-bearing idealisation individually: radiative recombination, mobility, series resistance to 5~$\Omega$cm$^2$, contacts, affinities, optics, and density of states (the ceiling-vs-record comparison is reported both ways because it flips with the density-of-states treatment). A pilot guardrailed-reporting checklist is proposed.

\textbf{Novelty and importance.} First, joint affinity mapping replaces one-at-a-time optimisation and moves the honest ceiling from 23.50\% (surveyed affinities) to 24.95\%. Second, the envelope is only as good as its prior, so the prior ranges and their measurement anchors are published in the main text, not the supplement. Third, scripted resistance sweeps to 5~$\Omega$cm$^2$ bound the six-point fill-factor gap against measured devices as roughly half resistive, half recombination-driven. The defended number is the 20.82\% routine-defect result, not the ceiling.

\textbf{Relation to reference [11].} This manuscript is intended as the electrical-loss extension and companion to Ref.~[11] (S.~Guruprasad et al., \textit{Solar Energy} 305 (2026) 114251), which implemented the wavelength-dependent (optical) half of the realistic-simulation protocol of Saidarsan et al.~[10]; the present work implements the electrical half on one fully worked stack. We believe publishing the two halves side by side in the same journal best serves the community.

\textbf{Fit to the journal.} As the official journal of the International Solar Energy Society, \textit{Solar Energy} is the natural home for a reproducibility study that disciplines photovoltaic performance modelling; the companion optical half [11] already appears in this journal.

\textbf{Statements.} This work is original, is not under consideration elsewhere, and has been approved by all authors. The authors declare no competing interests. All SCAPS results come from archived simulation receipts in the public repository (cited in the Data availability statement), so every number is reproducible. AI assistance was used for editing and consistency checking only, as declared in the manuscript; no results were generated by AI. Figures use the authors' own simulation outputs exclusively.

\textbf{Suggested reviewers.}
\begin{enumerate}
\item Dr.\ Thomas Unold, Helmholtz-Zentrum Berlin (perovskite doping density and device physics).
\item Dr.\ Joseph Luther, National Renewable Energy Laboratory (CsPbI3 quantum-dot and thin-film photovoltaics).
\item Dr.\ Faiz Arith, Universiti Teknikal Malaysia Melaka (SCAPS-1D perovskite device modelling).
\end{enumerate}
Contact details will be entered in Editorial Manager; the editors are of course free to choose other reviewers.

Thank you for your consideration. We look forward to your response.

\vspace{2em}
\noindent Sincerely,\\
Hussain Touhid Siddiquee (corresponding author)
\end{document}
